\documentclass[twocolumn]{article}
\usepackage[T1]{fontenc}
\usepackage[utf8]{inputenc}
\usepackage[dvips]{graphicx}
\begin{document}

\title{The Semmelweis effect}

\author{Roel Pompen}
\maketitle

\section*{\includegraphics{./TODO/boek/roel.eps}}
\begin{abstract}
Recent contradictions in the figures regarding the COVID-19 epidemic are beginning to change
our look on this crisis. In order to allow us to interpret these findings and get a
better view of where we are I'm delving into medical history. I hope to provide 
some reassurance for the next confrontation with figures that are simply incomplete and 
therefore do not allow any meaningful interpretation. Not the finest form of
reassurance, but it'll have to do.
\end{abstract}

\section*{Hygiene in the hospital}

In 1847 the general hospital in Vienna had two maternity clinics. One run by nurses
and one by doctors who at the time usually were surgeons also, as there wasn't that
much specialization going on. The clinic run by doctors had triple the lethalities that
the other ward had and that had been the case for roughly a century. (The hospital
was founded around 1700)

In 2 months time the death rate at the doctor's ward had plummeted to less than
that of the nurse's ward and after a year to almost zero. This phenomenal result was
brought on by the action of Ignaz Semmelweis who laid out a way of working based on
hygiene. And this before the work of famous people like Louis Pasteur and 
Joseph Lister (after whom Listerine mouth wash is named), who came up with theories 
of germs causing disease.

\subsection*{The figures}

In the clinic run by doctors the annual mortality rate was 10\%. The clinic run by nurses
had a rate of 4\%. The reason was that surgeons who taught anatomy in the basement making
use of cadavers, the stench of putrefaction dominating the room, went upstairs without washing
their hands and bringing their tools along, to come help in the maternity clinic.

Ignaz Semmelweis proposed to keep the instruments that were used for cadavers in the
basement and to clean tools used elsewhere in between uses. Furthermore doctors should wash
their hands in chlorinated lime when moving between rooms. This resulted in a drop in 
lethalities to practically 0 within 14 months.

\subsection*{Academic embarrassment}

Semmelweis wasn't seen as a hero, since the hospital was retroactively accused of being
responsible for a lot of deaths due to filthy behaviour. Keep in mind that hygiene was an
unfamiliar concept at the time, but insights change quickly. Everyone had lost relatives
from puerperal fever. Semmelweis was therefore socially exiled. Although accounts differ,
he was committed to a psychiatric hospital, possibly due to neurosyphilis or depression.
Fact is he died alone in a 19th century psych ward and his contemporaries wrote horrible
things about him. He allegedly died from gangrene after being beaten by guards.

\section*{The Semmelweis effect}

This brings us to the title of this paper. The Semmelweis effect is what I call the concept
that when a scientifically embarrassingly simple improvement leads to radical change in 
scientific prosperity, the instigators will become heretics, or progression is stifled until
the scientific establishment is either retired or dead.

I'll now try and illustrate that the establishment acts out of commercial interests and 
scientific delusion and that the time has arrived for our established ``experts'' to
clear the media stage, preferably before their retirement; A positive consequence of
the Corona debacle.

By the way this Semmelweis effect does not only appear in the medical sciences, it
appears in disciplines where you wouldn't expect it,  like mathematics and computer
science. The insight that the medical sciences are over represented can be explained by the
fact that physicians are used to do as taught in order to save lives. There is little room 
for doubting the subject matter. At the other end of the scientific spectrum we find the
theoretical physicist who, in order to function properly, is prepared to challenge everything
that was taught!

\section*{The relation with COVID-19}

People are slowly beginning to understand that we want to be advised by experts who have no
ties with the commercial sector. An epidemiologist has no business in the free market.
We know from epidemiology that a lockdown cannot be effective, something that is starting
to become clear from recent figures. Scientists are surprised that there are more people 
with immunity than assumed. 

As suggested earlier I'm not going to mention a lot of figures. After all these will change
over time. What I will do is present projections that'll be confirmed or discredited within
2 years or so. Obviously I'm writing  this article because I expect my predictions to hold.

\subsection*{Epidemiological aspects}

Areas with and without lockdown  seem to get comparable results and this was to be expected.
Assuming a lockdown effectiveness of 95\% (too high according to my estimates), a basic
reproduction number of 2 and a generation interval of 5.2 days (according to current
estimates), will delay the epidemic by 22 days.

Furthermore the demographic and geographic information becomes useless as we need to locate
the continuations of the epidemic. Basically all that you're achieving is two weeks
lockdown extension and 10 days of fighting the epidemic with your eyes closed.

\subsection*{Counterproductive}

The draconian countermeasures being taken as initiative of businesses and civilians usually
are counterproductive because these strategies require medical skills. You can't expect
people to understand and obey the protocols regarding gloves, for example.

And gloves are dangerous. The skin is our first barrier against infections by secreting
substances like defensins, lactic acid and fatty acids, which neutralize microbes.  On the
other hand it's the skin temperature which accelerates the multiplication of microbes.

By wearing gloves you bar the remedial aspects of the skin and retain the body
temperature.  Now you can breed you own microbes and administer them to your bloodstream 
as close to the brain as possible. In other words: A person touches his or her own face 320 
times per day. The face has mucous membranes at the mouth, eyes and nose that give access
to the bloodstream. 

If instead of wearing gloves you're a fan of using disinfecting gel five times a day you 
compromise defensin production and the skin on your hands loses much of its natural barrier 
against bacteria.

\subsection*{Damaging civil initiatives}
The fact that landlord of condos (at least here in Belgium) are taking the
initiative to shut down the elevators so that everyone has to touch the same
stair's handrail doesn't help either. I estimate that at least 50\% of the people
above 45 years of age can't safely walk the stairs with a bag of groceries without
using the rail with the other hand. In supermarkets in Belgium taking a trolley 
is mandatory because it's a good idea for people to all touch the same set of objects
\ldots Obviously they should ask for people if what they needs fits in a single bag
and if so supply one immediately.\footnote{As an aside: In Europe the concept 
of people bagging your groceries is unheard of and you even pay extra for the bag.}
And it doesn't stop there. 

People will claim that these measures are ok, because disinfectants are used constantly,
but that is reversed reasoning; Let's not collectively touch the same shopping carts
hoping that these are cleansed with diligence. The disinfecting strategy solves a problem
you wouldn't be having in the first place if you were allowed to forego the trolley when 
buying just a single bread!

\subsection*{Super-infection}

The reason this story about bacterial infection is relevant is that people with
more than one serious infection (a super-infection) fall inside the risk group. In 
other words if you get a bacterial infection and corona comes along your immune system 
has a tougher challenge to keep you alive.

\subsection*{Commercial interests}
We know since the work of the grandfather of epidemiology (William Farr, most
important publications around 1870) that an epidemic starts with an exponential growth
and flattens when reproduction becomes impossible. This is when too many people are
already immune.

Virologists are trying to change this curve using tests and medications, but it is known
from a mathematical point of view one can only have limited influence on this curve.
Venture capitalists are not aware of this and voila, a market is born.

It's these experts that constitute the establishment that should start clearing
the scene. We should be advised by independent specialists who base their ideas on
uncompromised science. Not the science of the plagiarism generation. After all that
science is mostly based on result oriented research.

\subsection*{Opportunism}

In these times people from big pharma are bending over backwards to generate funding
for their vaccine, which will be ready in two years after the epidemic. And which \emph{might} 
actually work with the next mutation. In the Netherlands the government has issued money
for having corona apps developed. Motivation and design goals unknown. These apps will 
violate privacy and work with facial recognition. While these apps are in development
we are now thinking about making masks mandatory, making these apps completely worthless.
At least we're burning money.

\subsection*{Ill advice and figures}

\begin{table}
\noindent \begin{centering}
\begin{tabular}{|c|c|c|}
\hline 
\multicolumn{3}{|c|}{people / \% people tested}\tabularnewline
\hline 
\textbf{seronegative} &  & \textbf{seropositive}\tabularnewline
\hline 
people / \% &  & people / \%\tabularnewline
\hline 
\multicolumn{3}{|c|}{\textbf{outside risk group}}\tabularnewline
\hline 
\textbf{hospitalized} & \textbf{deceased} & \textbf{cured}\tabularnewline
\hline 
people / \% & people / \% & people / \%\tabularnewline
\hline 
\multicolumn{3}{|c|}{\textbf{inside risk group}}\tabularnewline
\hline 
\textbf{hospitalized} & \textbf{deceased} & \textbf{cured}\tabularnewline
\hline 
people / \% & people / \% & people / \%\tabularnewline
\hline 
\end{tabular}
\par\end{centering}
\caption{Snapshot of an epidemic}
\end{table}

An epidemiologist who contributes to a TV show that comments on figures that are incomplete 
has the moral and academic obligation to defy the presentation of those incomplete facts and
possibly even bring his or her own figures along. A suggestion for a complete set of data
can be found in the supplied table. Obviously such a set of numbers should be accompanied by
a description of the risk group, to overcome the problem of people being classified as having died of
corona while there is actually an underlying condition to which the death should be attributed.
This is why we see the number of lethal heart attacks dropping. A part of these people now
count as corona victims!

\section*{Conclusion}

An expert who doesn't support the presentation of a complete set of fact creates uncertainty.
Uncertainty creates panic. Panic creates pressure on politicians. Politicians deviate from
the initially planned epidemiological tactics (``Draconian measures''). Civilians take initiatives
that match the vigor of the government. And voila.

It's not like people go en masse to the hospital to die. People are dying because people go
en masse to the hospital\ldots or other places. Scared people cluster. At the drugstore,
church (in the Netherlands religious congregations up to 30 people are allowed.) Lockdown parties.
Recreational parks.

What is happening? We don't know anymore who the experts are. Governments should have 
independent (non-affiliated) epidemiologists who determine the strategy to follow and 
present this to the public at least twice a week, using complete figures. This way they 
take away opportunities for other experts to present themselves in the media, with their
convenient omissions. This breaks the current vicious cycle  of panic and political
adjustments so that we can assure this never happens again.


\end{document}
